\section{Finite cycles, product identities and rank inequalities}

\begin{lemma}[Finite-cycle coboundary]
    \label{lem:finite_cycle_coboundary}
    \lean{Fin.exists_succ_coboundary_of_prod_eq_one,
        Fin.exists_cyclic_coboundary_of_prod_eq_one,
        Fin.exists_complex_unit_cyclic_coboundary_of_prod_eq_one}
    \leanok
    Let $G$ be a commutative group, and let
    $\kappa_0,\ldots,\kappa_n\in G$. If
    $\prod_{k=0}^{n}\kappa_k=1$, then there are
    $\varphi_0,\ldots,\varphi_n\in G$ such that
    $\kappa_k=\varphi_k\varphi_{k+1}^{-1}$ for every $k<n$, and
    $\kappa_n=\varphi_n\varphi_0^{-1}$.
    Equivalently, on a nonempty finite cyclic index set, a product-one family
    satisfies $\kappa_k=\varphi_k\varphi_{k+1}^{-1}$ for every cyclic edge.
    In particular, a product-one family of complex unit phases admits such a
    choice with all $\varphi_k$ again of unit modulus.
\end{lemma}

\begin{proof}
    \leanok
    Set $\varphi_k=(\kappa_0\cdots\kappa_{k-1})^{-1}$, with the empty product
    equal to $1$. For $k<n$,
    \begin{align}
        \varphi_k\varphi_{k+1}^{-1}
         & =
        (\kappa_0\cdots\kappa_{k-1})^{-1}
        (\kappa_0\cdots\kappa_k)
        =\kappa_k. \notag
    \end{align}
    Separating the final factor in the product-one hypothesis gives
    $(\kappa_0\cdots\kappa_{n-1})\kappa_n=1$, which gives
    \begin{align}
        \varphi_n\varphi_0^{-1}
        =
        (\kappa_0\cdots\kappa_{n-1})^{-1}
        =
        \kappa_n. \notag
    \end{align}
    For unit phases the same construction runs in the circle group.
\end{proof}

\begin{lemma}[Coordinates for a cyclic step]
    \label{thm:periodic_step_coordinates}
    \lean{Fin.stepOrbit_bijective, Fin.stepOrbitEquiv, Fin.stepOrbitEquiv_add_one}
    \leanok
    Let $m>0$, $p\geq0$, and $r=\gcd(m,p)$. Every residue $u$ modulo $m$
    has a unique expression $u=\alpha+pk\pmod m$, with
    $0\leq\alpha<r$ and $0\leq k<m/r$.
    In these coordinates from
    \cite[proof of Theorem~4.1]{DeLasCuevas2017Irreducible}, adding $p$
    preserves $\alpha$ and increments $k$ by one modulo $m/r$.
\end{lemma}

\begin{proof}
    \leanok
    If two such expressions agree modulo $m$, reduction modulo $r$
    identifies their representatives $\alpha$. Cancelling that common
    representative gives $pk=pk'\pmod m$, hence $k=k'\pmod{m/r}$.
    Their bounds imply $k=k'$. There are $r(m/r)=m$ pairs of coordinates,
    so this injective map to the $m$ residues is also surjective.
    Since $m$ divides $p(m/r)$, reducing $k+1$ modulo $m/r$ does not
    change $\alpha+p(k+1)$ modulo $m$.
\end{proof}

\begin{lemma}[Orbit roots as consecutive products]
    \label{thm:periodic_orbit_phase_products}
    \lean{Fin.exists_stepOrbit_phases}
    \leanok
    \uses{thm:periodic_step_coordinates}
    Let $G$ be an abelian group, $m>0$, $p\ge0$, $r=\gcd(m,p)$,
    and $q=m/r$. Suppose $c_\alpha\in G$ satisfies $c_\alpha^q=1$
    for every $0\le\alpha<r$. There are $b_u\in G$, indexed modulo
    $m$, such that, writing $u=\alpha_u+pk_u$ for the coordinates of
    Lemma~\ref{thm:periodic_step_coordinates},
    \begin{align}
        \prod_{j=0}^{p-1}b_{u+j} & = c_{\alpha_u}.
        \label{eq:periodic_orbit_phase_product}
    \end{align}
    For $G=U(1)$, these are the local phases in
    \cite[proof of Theorem~4.1]{DeLasCuevas2017Irreducible}, where the
    coordinate $k_u$ is now advanced modulo $q$.
\end{lemma}

\begin{proof}
    \leanok
    \uses{thm:periodic_step_coordinates}
    Put $f_u=c_{\alpha_u}^{k_u}$ and $b_u=f_{u+1}/f_u$.
    The coordinate shift and $c_\alpha^q=1$ give
    $f_{u+p}=c_{\alpha_u}f_u$, including when $k_u=q-1$.
    The product telescopes to $f_{u+p}/f_u=c_{\alpha_u}$,
    proving~(\ref{eq:periodic_orbit_phase_product}).
\end{proof}

\begin{lemma}[Multiplication of an orthogonal resolution]
    \label{thm:periodic_orthogonal_resolution_multiplication}
    \lean{orthogonalProjection_mul_eq_ite_of_sum_eq_one}
    \leanok
    If orthogonal projections $P_u$ satisfy $\sum_u P_u=\Id$, then
    $P_uP_v=\delta_{uv}P_u$.
\end{lemma}

\begin{proof}
    \leanok
    Each projection is idempotent, $P_u^2=P_u$, and distinct projections in an
    orthogonal resolution have orthogonal ranges, $P_uP_v=0$ for $u\ne v$.
\end{proof}

\begin{lemma}[Eventually constant weighted cycle sums]
    \label{thm:beigi_eventually_constant_sector_graph}
    \lean{FiniteWeightedDigraph.cyclicEdgeWeight_eq_one,
        FiniteWeightedDigraph.cycle_eq_constant,
        FiniteWeightedDigraph.cycleWeightSum_eq_card_loop}
    \leanok
    Let $G$ be a finite directed graph whose edge weights are non-negative
    integers, and let
    \begin{align}
        s_N
         & =
        \sum_{(a_0,\ldots,a_{N-1})}
        \prod_{n=0}^{N-1}w(a_n,a_{n+1}) \notag
    \end{align}
    be the weighted sum over ordered cyclic vertex sequences, indices modulo
    $N$. Suppose that $s_N$ is constant for all sufficiently large $N$. By
    the comparison in \cite[Section~IV]{Beigi2012Classification}, every edge
    that belongs to a directed cycle of positive weight then has weight one,
    every such cycle is a repetition of a single loop, and, for every $N>0$,
    $s_N=\#\{a:w(a,a)>0\}$.
\end{lemma}

\begin{proof}
    \leanok
    For a sufficiently large $N$, the sums at lengths $N$,
    $N+1$, and $N(N+1)$ share their common eventual value. Repeating each
    ordered $N$-cycle $N+1$ times gives an injection into the cycles of length
    $N(N+1)$, while repeating each ordered $(N+1)$-cycle $N$ times gives a
    second injection. Comparison of the three weighted sums shows first that
    every positive cyclic weight is one. Both injections must then be
    exhaustive. Hence a cycle at the common length has periods $N$ and
    $N+1$; comparison of adjacent entries gives period one. Repeating an
    arbitrary shorter cycle to a sufficiently large length proves that it is
    constant. The only surviving ordered cycles at any positive length are
    therefore the constant words supported on positive loops, each of weight
    one.
\end{proof}

\begin{lemma}[Factor scalars from a product tensor identity]
    \label{thm:periodic_product_tensor_phase_extraction}
    \lean{PiTensorProduct.exists_kappa_of_tprod_eq_smul}
    \leanok
    Let $z\in\C$ be nonzero. For each position $v=0,\ldots,m-1$, let
    $\{A_v^i\}_i$ and $\{B_v^i\}_i$ be families of complex matrices with the
    same rectangular size. Suppose that for every choice of letters
    $\sigma=(\sigma_v)_v$,
    \begin{align}
        \bigotimes_{v=0}^{m-1} A_v^{\sigma_v}
         & =
        z\,\bigotimes_{v=0}^{m-1} B_v^{\sigma_v}.
        \notag
    \end{align}
    Suppose also that for each $v$ there is a reference letter $i_v$ and a
    matrix entry $(r_v,c_v)$ with
    $(B_v^{i_v})_{r_v,c_v}\ne 0$. Then there are scalars
    $\kappa_v\in\C$ such that
    \begin{align}
        \prod_{v=0}^{m-1}\kappa_v & =z,
        \notag                                      \\
        A_v^i                     & =\kappa_v B_v^i
        \notag
    \end{align}
    for every position $v$ and every letter $i$. This is the extraction step
    of \cite[Appendix~A, lines~1072--1080]{DeLasCuevas2017Irreducible}.
\end{lemma}

\begin{proof}
    \leanok
    For each position choose the coordinate functional at the nonzero
    reference entry and normalize it so that it takes the value $1$ on
    $B_v^{i_v}$. Applying the tensor product of these functionals to the
    reference identity gives $\prod_v\kappa_v=z$. To prove the
    proportionality at position $v$, replace only the $v$-th functional by the
    coordinate functional at an arbitrary entry $(r,c)$ and apply the
    identity with $\sigma_v=i$ and all other letters fixed at their
    references. The product over the other positions is nonzero because its
    product with $\kappa_v$ is $z$, so cancellation gives
    $(A_v^i)_{r,c}=(\kappa_v B_v^i)_{r,c}$ for every entry.
\end{proof}

\begin{lemma}[Product-one factor scalars after choosing a root]
    \label{thm:periodic_product_one_phase_extraction}
    \lean{PiTensorProduct.exists_kappa_product_one_of_tprod_eq_root_smul}
    \leanok
    \uses{thm:periodic_product_tensor_phase_extraction}
    In the setting of
    Lemma~\ref{thm:periodic_product_tensor_phase_extraction}, let $m>0$ and
    choose $\xi\in\C$ with $\xi^m=z$. Then there are scalars $\kappa_v\in\C$
    with $\prod_v\kappa_v=1$ and $A_v^i=\kappa_v\xi\,B_v^i$ for every
    position $v$ and every letter $i$.
\end{lemma}

\begin{proof}
    \leanok
    Lemma~\ref{thm:periodic_product_tensor_phase_extraction} yields
    scalars $\mu_v$ with $\prod_v\mu_v=z$ and $A_v^i=\mu_vB_v^i$. Since
    $m>0$, the identities $\xi^m=z$ and $z\ne0$ imply $\xi\ne0$, so the
    ratios $\kappa_v=\mu_v/\xi$ satisfy
    $\prod_v\kappa_v=(\prod_v\mu_v)/\xi^m=z/z=1$ and
    $A_v^i=(\kappa_v\xi)B_v^i$.
\end{proof}

\begin{definition}[Common buffer length of a finite family]
    \label{def:common_buffer_length}
    \lean{commonBufferLength, commonBufferLength_pos,
        le_commonBufferLength_add_one, commonBufferLength_le}
    \leanok
    For a finite family of natural numbers $(N_i)_{i\in I}$, the
    \emph{common buffer length} is
    \begin{align}
        L & =\max\Bigl(1,\max_{i\in I}(N_i-1)\Bigr),
        \label{eq:common_buffer_length}
    \end{align}
    with $N_i-1$ read as $0$ when $N_i=0$.  It is positive, satisfies
    $N_i\leq L+1$ for every $i\in I$, and is the least positive integer with
    this property.
\end{definition}

\begin{lemma}[Choi rank and the Kraus span]
    \label{thm:periodic_choi_rank_kraus_span}
    \lean{Channel.choiRank_mapLM_eq_finrank_span,
        Channel.choiRank_mapLM_eq_finrank_wordSpan_one}
    \leanok
    Let $\mathcal F(X)=\sum_a A^aX(A^a)^*$ on $\MN{D}$, and let
    $S=\spn_{\C}\{A^a\}_a$. The Choi rank of $\mathcal F$ is
    $\dim_{\C}S$, including when $D=0$.
\end{lemma}

\begin{proof}
    \leanok
    For $D>0$, let $C$ have columns equal to the vectorized Kraus operators,
    multiplied by the common nonzero normalization of the Choi matrix. That
    matrix is $CC^*$, whose rank is the rank of $C$. Vectorization preserves
    the dimension of the Kraus span. For $D=0$, both ranks are zero.
\end{proof}

\begin{lemma}[Rank under an injective positive map]
    \label{thm:periodic_injective_positive_rank}
    \lean{Matrix.rank_le_rank_map_of_injective_positive,
        Matrix.rank_le_rank_map_of_injective_positive_fintype}
    \leanok
    Let $\Phi:\MN{D}\to\MN{D}$ be an injective positive linear map.
    If $A\ge0$, then $\operatorname{rank}A\le\operatorname{rank}\Phi(A)$.
\end{lemma}

\begin{proof}
    \leanok
    The positive matrices bounded above by a positive multiple of $A$
    span the matrix algebra on its support. Positivity carries this span
    into the analogous span for $\Phi(A)$. Their complex dimensions are
    $(\operatorname{rank}A)^2$ and $(\operatorname{rank}\Phi(A))^2$,
    respectively, so injectivity gives the inequality.
\end{proof}

\begin{lemma}[Choi rank under injective composition]
    \label{thm:periodic_injective_composition_choi_rank}
    \lean{Channel.choiRank_le_comp_of_injective,
        Channel.choiRank_le_pow_of_injective}
    \leanok
    Let $\mathcal F$ and $\mathcal G$ be completely positive maps on
    $\MN{D}$, and suppose $\mathcal F$ is injective. Then
    $\operatorname{rank}C_{\mathcal G}\le
        \operatorname{rank}C_{\mathcal F\circ\mathcal G}$.
    Consequently $\operatorname{rank}C_{\mathcal F}\le
        \operatorname{rank}C_{\mathcal F^p}$ for every $p>0$.
\end{lemma}

\begin{proof}
    \leanok
    \uses{thm:periodic_injective_positive_rank}
    Apply Lemma~\ref{thm:periodic_injective_positive_rank} to the
    injective positive map $\mathcal F\otimes\mathrm{id}$ and the
    positive Choi matrix $C_{\mathcal G}$. The image is
    $C_{\mathcal F\circ\mathcal G}$. Iteration gives the power inequality.
\end{proof}
